%! TeX root = ../phd-thesis.tex
%! Author = davidedomini
\chapter{Reinforcement Learning}
\label{chap:reinforcement-learning}
\minitoc

\gls{RL} is a learning paradigm in which an agent discovers
 a behavioral policy through direct interaction with an environment,
 guided only by a scalar reward signal.
%
This chapter introduces the theoretical foundations needed to follow the contributions
 developed in subsequent chapters.
%
\Cref{sec:rl-framework} describes the general \gls{RL} framework at a conceptual level,
 introducing the key actors and the feedback loop that drives learning.
%
\Cref{sec:rl-single-agent} then formalises the single-agent setting through the \gls{MDP}
 and presents \gls{DQL} as the function-approximation algorithm used throughout this work.
%
\Cref{sec:rl-multi-agent} extends the treatment to cooperative multi-agent systems,
 introducing the \gls{MARL} formalism, the SwarMDP abstraction for homogeneous swarms,
 and the \gls{CTDE}/\gls{DTDE} paradigms.
%
\Cref{sec:rl-scalability} addresses the scalability bottleneck that arises
 when the population grows large,
 presenting \glspl{GNN} as a principled mechanism for size-agnostic,
 permutation-invariant policies and zero-shot transfer.
%
Finally, \Cref{sec:rl-decentralized} examines fully decentralized training strategies
 in which agents improve their policies through peer-to-peer exchanges
 rather than through a central coordinator.


%=============================================================================
\section{The Reinforcement Learning Framework}
\label{sec:rl-framework}
%=============================================================================

The defining characteristic of \gls{RL} is that the learning signal comes not from
 labelled examples but from the consequences of the agent's own actions~\cite{Sutton1998}.
%
Two entities interact: an \emph{agent} and an \emph{environment}.
%
At each discrete time step the agent perceives a representation of the environment's
 current \emph{state}, chooses an \emph{action} from the set of available alternatives,
 and in response the environment transitions to a new state and emits a scalar \emph{reward}.
%
This loop repeats indefinitely, or until a terminal condition is reached,
 producing a trajectory of states, actions, and rewards.

The agent's goal is to find a \emph{policy} (i.e., a mapping from states to actions) 
 that maximizes the cumulative reward collected over time.
%
What makes the problem non-trivial is that reward is often delayed:
 an action taken now may not produce observable benefit until many steps later,
 forcing the agent to reason about long-horizon consequences rather than greedy immediate gains.
%
This is the \emph{credit assignment} problem, i.e., attributing future outcomes
 to the earlier decisions responsible for them.

A further challenge is the \emph{exploration--exploitation trade-off}.
%
To collect reward the agent must act on its current knowledge of the environment
 (\emph{exploit}), but to improve that knowledge it must occasionally try actions
 it has not taken before (\emph{explore}).
%
Resolving this tension efficiently is a central concern in \gls{RL} algorithm design.

\begin{figure}[htbp]
  \centering
  \fbox{\parbox[c][4.0cm][c]{0.88\linewidth}{\centering
    \textbf{Image placeholder}\\[4pt]
    The agent--environment interaction loop:
    at each step the agent observes state $s_t$, selects action $a_t$,
    and receives reward $r_{t+1}$ together with the next state $s_{t+1}$.
    Arrows show the cyclic flow of perception, decision, and feedback.
  }}
  \caption{The agent--environment interaction loop in \gls{RL}.}
  \label{fig:rl-loop}
\end{figure}

Unlike supervised learning, which assumes a fixed labelled dataset
 independent of the learner's behavior,
 \gls{RL} faces a \emph{non-stationary data distribution}:
 the states the agent encounters depend on the policy it follows,
 which in turn changes as learning progresses.
%
This coupling between behavior and data collection is what makes \gls{RL} both
 more general than supervised approaches and more challenging to optimize.
%
The following sections introduce the mathematical framework
 that makes these intuitions precise.

%=============================================================================
\section{Single-Agent Reinforcement Learning}
\label{sec:rl-single-agent}
%=============================================================================

The standard formalization of a sequential decision-making problem is 
 the \glsfull{MDP}~\cite{Sutton1998},
 defined as a tuple $(\mathcal{S}, \mathcal{A}, P, R, \gamma)$.
%
$\mathcal{S}$ is the state space and $\mathcal{A}$ is the action space.
%
$P : \mathcal{S} \times \mathcal{A} \times \mathcal{S} \to [0,1]$ is the transition
 probability kernel, encoding how the environment evolves in response to the agent's actions.
%
$R : \mathcal{S} \times \mathcal{A} \to \mathbb{R}$ is the reward function,
 and $\gamma \in [0,1)$ is the discount factor that controls the relative importance
 of immediate versus future rewards.
%
At each discrete time step $t$ the agent observes the current state $s_t \in \mathcal{S}$,
 selects an action $a_t \in \mathcal{A}$ according to a policy
 $\pi : \mathcal{S} \to \Delta(\mathcal{A})$, receives a scalar reward
 $r_{t+1} = R(s_t, a_t)$, and transitions to $s_{t+1} \sim P(\cdot \mid s_t, a_t)$.

The goal of the agent is to find a policy $\pi^*$ that maximizes the expected discounted return.
%
The central object for this purpose is the action-value function $Q^\pi(s, a)$,
 which quantifies the expected return obtained by taking action $a$ in state $s$
 and following $\pi$ thereafter:

\begin{equation}
  Q^\pi(s, a) = \mathbb{E}_\pi\!\left[\sum_{t=0}^{\infty} \gamma^t r_{t+1}
    \;\middle|\; s_0 = s,\, a_0 = a\right].
  \label{eq:q-value}
\end{equation}

The optimal action-value function $Q^*(s,a) = \max_\pi Q^\pi(s,a)$ satisfies
 the Bellman optimality equation:

\begin{equation}
  Q^*(s, a) = \mathbb{E}_{s' \sim P}\!\left[r + \gamma \max_{a' \in \mathcal{A}}
    Q^*(s', a') \;\middle|\; s, a\right],
  \label{eq:bellman}
\end{equation}

from which the optimal policy is recovered greedily as $\pi^*(s) = \arg\max_{a} Q^*(s,a)$.

For large or continuous state spaces, tabular representations of $Q^*$ are infeasible.
%
\glsfull{DQL}~\cite{mnih2015humanlevel} approximates the action-value function with
 a neural network $Q(s, a; \theta)$ parameterized by weights $\theta$, 
 trained by minimising a temporal-difference loss with respect to a periodically frozen
 \emph{target network} $Q^- = Q(\cdot;\theta^-)$:

\begin{equation}
  \mathcal{L}(\theta) = \mathbb{E}_{\langle s,a,r,s'\rangle \sim \mathcal{D}}\!\left[
    \!\left(r + \gamma \max_{a'} Q(s', a';\theta^-) - Q(s, a;\theta)\right)^{\!2}\right],
  \label{eq:dql-loss}
\end{equation}

where $\mathcal{D}$ is a \emph{replay buffer} of past transitions.
%
Two mechanisms are responsible for stabilising the optimization.
%
The target network, whose weights $\theta^-$ are copied from $\theta$ only periodically,
 provides a fixed regression target and prevents the gradient from chasing a moving signal.
%
Experience replay breaks the temporal correlation of online data by sampling mini-batches
 uniformly from $\mathcal{D}$, making the training distribution closer to i.i.d.

\begin{figure}[htbp]
  \centering
  \fbox{\parbox[c][4.2cm][c]{0.88\linewidth}{\centering
    \textbf{Image placeholder}\\[4pt]
    Architecture of a Deep Q-Network: input state $s$,
    hidden layers, output layer emitting $Q(s,a;\theta)$ for each action $a \in \mathcal{A}$;
    arrows showing TD-loss gradient flow and the periodic copy $\theta \to \theta^-$
    to the target network.
  }}
  \caption{Deep Q-Network architecture.}
  \label{fig:dqn-architecture}
\end{figure}

%=============================================================================
\section{Multi-Agent Reinforcement Learning}
\label{sec:rl-multi-agent}
%=============================================================================

When several agents interact within a shared environment their joint decisions determine
 the trajectory of the system, and a single \gls{MDP} is no longer sufficient to model it.
%
\gls{MARL}~\cite{DBLP:journals/smca/BusonuiBDS08,DBLP:journals/corr/abs-1911-10635}
 is the extension of \gls{RL} to this multi-agent setting,
 and the nature of the agents' objectives gives rise to three broad classes of problems.
%
In \emph{competitive} settings (the paradigmatic case being two-player zero-sum games),
 the gain of one agent is the loss of another, so agents must learn to act adversarially.
%
\emph{Mixed} or general-sum settings allow both competitive and cooperative incentives
 to coexist simultaneously, as in market simulations or mixed-motive social dilemmas.
%
\emph{Cooperative} settings, finally, require all agents to maximize a shared objective:
 every agent receives the same team reward, and the goal is to find a joint policy that
 maximizes the expected collective return.
%
This thesis focuses exclusively on cooperative \gls{MARL},
 which is the natural formulation for collective systems where agents must coordinate
 to accomplish tasks that no single agent could achieve alone.
%
In this setting the standard formalism augments the \gls{MDP} with a communication structure.
%
The setting is described by the tuple $(\mathcal{G}, \mathcal{S}, \mathcal{A},
 \mathcal{O}, P, R, \gamma)$, where $\mathcal{G} = (\mathcal{N}, \mathcal{E})$
 is a graph whose nodes are agents and whose edges represent communication links,
 $\mathcal{O}$ is the per-agent observation space,
 and the remaining components generalize those of the single-agent \gls{MDP}.
%
Because agents typically cannot observe the full global state,
 each agent $i \in \mathcal{N}$ receives only a local observation $o_i^t \in \mathcal{O}$,
 a setting formally captured by the Dec-POMDP model~\cite{Oliehoek2016}.

For large \emph{homogeneous} swarms, where every agent runs the same policy prototype,
 the SwarMDP model~\cite{DBLP:conf/atal/SosicKZK17} provides a specialized formalism,
 adopted and extended in subsequent work~\cite{DBLP:journals/ai/AguzziDVV26,DBLP:journals/fgcs/FarabegoliDAV26}.
%
An agent prototype is a tuple $\mathbf{a} = (\mathcal{S}, \mathcal{O}, \mathcal{A}, \rho, \pi)$
 comprising local state, observation, and action spaces,
 an internal update function $\rho$, and a policy $\pi$.
%
The environment is characterized by the population size $N$,
 a global transition function $T$, and an observation function $Y$.
%
The collective dynamics unfold as a functional sequence applied element-wise
 across all $N$ agents:

\begin{equation}
  \mathcal{S}^N \xrightarrow{Y} \mathcal{O}^N \xrightarrow{\pi} \mathcal{A}^N \xrightarrow{T} \mathcal{S}^N.
  \label{eq:swarmdp-cycle}
\end{equation}

Stacking individual observations, rewards, and actions into global vectors
 $\mathbf{o}^t = (o_1^t, \ldots, o_N^t)$, $\mathbf{r}^t$, $\mathbf{a}^t$,
 the time-varying neighborhood of agent $i$ is $\mathcal{N}_t(i) \subseteq \{1, \ldots, N\}$,
 inducing the communication graph:

\begin{equation}
  G_t = (\mathcal{V},\, \mathcal{E}_t), \qquad
  \mathcal{V} = \{1,\ldots,N\}, \qquad
  \mathcal{E}_t = \{(i,j) \mid i \in \mathcal{V},\, j \in \mathcal{N}_t(i)\}.
  \label{eq:swarmdp-graph}
\end{equation}

The collective objective is to maximize the expected cumulative team reward:

\begin{equation}
  G = \sum_{t=0}^{H-1} \sum_{i=1}^{N} r_i^{t+1}, \qquad
  \pi^* = \arg\max_\pi\, \mathbb{E}_\pi[G].
  \label{eq:swarmdp-objective}
\end{equation}

Each agent $i$ maintains its own action-value estimate conditioned on its local observation:

\begin{equation}
  Q_i^{\pi,t}(o_i^t, a_i^t) = \mathbb{E}_\pi\!\left[
    \sum_{k=t}^{T} \gamma^{k-t} r_i^k \;\middle|\; o_i^t, a_i^t\right].
  \label{eq:per-agent-q}
\end{equation}

Training follows the same \gls{DQL} paradigm as in the single-agent case,
 substituting the per-agent quantities into the loss of \Cref{eq:dql-loss}.

The literature distinguishes two orthogonal paradigms for cooperative \gls{MARL}~\cite{DBLP:conf/nips/LoweWTHAM17,DBLP:conf/icml/RashidSWFFW18}.
%
\glsfull{CTDE} grants a central learner access to the global state and all agents' observations
 during training, while at execution time each agent acts using only local information.
%
\glsfull{DTDE}, conversely, requires both training and execution to be carried out
 with local information only, without any central coordinator at any stage.
%
\gls{CTDE} typically yields higher-quality policies because the learner can observe
 the full joint state, but it requires a central coordinator and is therefore unsuitable
 for settings where agents must remain autonomous throughout.
%
\gls{DTDE} is more resilient to communication failures and easier to deploy on
 resource-constrained hardware, at the cost of a harder optimization problem.

\begin{figure}[htbp]
  \centering
  \fbox{\parbox[c][4.5cm][c]{0.88\linewidth}{\centering
    \textbf{Image placeholder}\\[4pt]
    Side-by-side comparison of \gls{CTDE} (left) and \gls{DTDE} (right).
    Left: a central critic receives global observations $\mathbf{o}^t$ and $G_t$ during training;
    individual actors use only $o_i^t$ at execution.
    Right: each agent trains and executes locally,
    exchanging only lightweight messages with direct neighbors.
  }}
  \caption{\gls{CTDE} vs.\ \gls{DTDE} paradigms in cooperative \gls{MARL}.}
  \label{fig:ctde-dtde}
\end{figure}

%=============================================================================
\section{Scalability in Multi-Agent Reinforcement Learning}
\label{sec:rl-scalability}
%=============================================================================

As the number of agents grows, cooperative \gls{MARL} faces two intertwined challenges.
%
The first is combinatorial: the joint action space $\mathcal{A}^N$ grows exponentially
 with $N$ (i.e., the number of agents), making centralized Q-value estimation impractical for large populations.
%
The second is architectural: a policy network whose input concatenates all agents'
 observations has a fixed input dimension, so it cannot be deployed on a population
 of different size from the one it was trained on.

One family of methods addresses the first challenge through \emph{value decomposition}.
%
Value Decomposition Networks~\cite{DBLP:conf/aaai/SunehagLGCZJLSL18} represent the
 joint action-value function as a sum of per-agent utility functions,
 factoring a complex joint optimization into $N$ independent local ones.
%
QMIX~\cite{DBLP:conf/icml/RashidSWFFW18} relaxes the additivity constraint to a
 monotonicity condition, allowing a mixing network to combine individual utilities
 in a more expressive way while still permitting a greedy joint action selection.
%
These approaches reduce the optimization complexity but still require centralized
 training and do not in themselves address the fixed-population limitation.

A complementary perspective comes from \emph{mean field approximation}~\cite{DBLP:conf/icml/YangLLZZWY18}.
%
Rather than modelling the interaction of agent $i$ with every other agent individually,
 mean field theory approximates the joint effect of the neighbors as a single
 mean action, reducing an $N$-body problem to a two-body problem between agent $i$
 and the aggregate of its neighborhood.
%
This allows the value function to be parameterized independently of $N$,
 but the approximation quality degrades when interactions are highly heterogeneous
 or the neighborhood distribution changes abruptly.

The most principled solution to the second challenge of architectural generality across
 population sizes comes from \glspl{GNN}~\cite{Scarselli2009,kipf2017semisupervised}.
%
A \gls{GNN} processes the communication graph $G_t$ through iterative message passing
 (\Cref{eq:gnn-message,eq:gnn-aggregation,eq:gnn-update}):
 each node aggregates information from its neighbors, updates its own embedding,
 and the operation is independent of the total number of nodes.
%
Graph Convolutional Reinforcement Learning~\cite{DBLP:conf/iclr/JiangDLSZ20}
 was an early demonstration of this idea, using graph convolutions to encode
 the local neighborhood of each agent into a relational representation
 before the Q-head.
%
Because the message-passing weights are shared across all nodes,
 the same network can process graphs of any size at test time,
 enabling \emph{zero-shot transfer}: a policy trained on $N$ agents
 generalizes directly to unseen population sizes without any fine-tuning,
 simply by running the same computation on a graph with more nodes.

\emph{Parameter sharing} is a further scalability lever orthogonal to the above~\cite{DBLP:conf/nips/LoweWTHAM17}.
%
When all agents are homogeneous, giving them identical network weights forces them
 to share representations and symmetrises the learning problem,
 effectively reducing the number of free parameters from $O(N \cdot |\theta|)$ to $O(|\theta|)$.
%
Combined with a \gls{GNN} backbone, parameter sharing and zero-shot transfer
 together make it possible to train on a small population and deploy on a much larger one,
 which is the key property required by swarm applications where
 $N$ may vary across deployments.

\begin{figure}[htbp]
  \centering
  \fbox{\parbox[c][5.0cm][c]{0.88\linewidth}{\centering
    \textbf{Image placeholder}\\[4pt]
    Comparison of scalability approaches in cooperative \gls{MARL}.
    Left: value decomposition (QMIX) --- centralized mixing network at training time.
    Centre: mean field --- each agent interacts with the mean action of its neighborhood.
    Right: GNN-based policy --- message passing over $G_t$, same weights for any $N$,
    zero-shot transfer to larger populations.
  }}
  \caption{Overview of scalability strategies in cooperative \gls{MARL}.}
  \label{fig:rl-scalability}
\end{figure}

%=============================================================================
\section{Decentralized Training in Multi-Agent Reinforcement Learning}
\label{sec:rl-decentralized}
%=============================================================================

A distinctive difficulty of \gls{MARL} that does not arise in the single-agent case
 is \emph{non-stationarity}: from the perspective of any single agent,
 the environment is non-stationary because the other agents are simultaneously
 updating their policies.
%
The simplest response is \emph{Independent Q-Learning}~\cite{tan1993multi},
 in which each agent trains its own Q-network while treating the other agents
 as part of a stationary environment.
%
Although this approach is easy to implement and scales trivially,
 it provides no convergence guarantees and performs poorly in strongly cooperative
 tasks where coordinated behavior is essential.

The standard alternative is \gls{CTDE}~\cite{DBLP:conf/nips/LoweWTHAM17,DBLP:conf/icml/RashidSWFFW18}
 (introduced in \Cref{sec:rl-multi-agent}): a central learner observes the global state
 during training and stabilizes the Q-estimates by conditioning them on information
 unavailable at execution.
%
While effective, this approach requires a central coordinator,
 a communication channel to aggregate all agents' data,
 and the ability to observe the full joint state: requirements that are difficult
 to satisfy in open, large-scale, or privacy-sensitive deployments.

Fully decentralized training removes the central coordinator entirely.
%
\citeauthor{DBLP:conf/icml/ZhangYBH18} formalise this setting through networked
 agents~\cite{DBLP:conf/icml/ZhangYBH18}:
 each agent communicates only with its direct neighbors on a fixed graph,
 performs local stochastic gradient steps on its own reward signal,
 and periodically averages its parameters with those neighbors via a
 consensus step borrowed from distributed optimization.
%
Under standard connectivity and stepsize conditions, this procedure converges
 to a neighborhood of the globally optimal cooperative policy,
 with a rate that depends on the spectral gap of the communication graph.

Communication of learned representations (rather than raw parameters or gradients) is
 explored by approaches such as CommNet~\cite{DBLP:conf/nips/SukhbaatarSF16} and
 DIAL~\cite{DBLP:conf/icml/FoersterAFW16}.
%
In these frameworks agents exchange continuous-valued messages differentiably,
 so the communication protocol is itself learned end-to-end alongside the policy.
%
CommNet uses a mean aggregation of hidden states to produce the message,
 whereas DIAL allows gradient flow through the communication channel during centralized
 training and then discretises the messages at execution.
%
Both methods demonstrate that allowing agents to learn \emph{what} to communicate,
 rather than prescribing the message content, yields qualitatively better coordination
 in tasks where the relevant information is task-specific and not known a priori.

The general principle underlying all decentralized training approaches is a
 trade-off between the quality of the learned policy and the independence
 of the learning process.
%
Methods that share more information (full states, all parameters, or learned messages) tend
 to converge faster and to higher-quality policies, but require stronger infrastructure
 and coordination assumptions.
%
Methods that restrict sharing to local neighborhoods sacrifice some performance
 for resilience, deployability, and privacy,
 making them the natural choice for large or resource-constrained systems.

\begin{figure}[htbp]
  \centering
  \fbox{\parbox[c][4.2cm][c]{0.88\linewidth}{\centering
    \textbf{Image placeholder}\\[4pt]
    Spectrum of training paradigms from fully centralized (left) to fully decentralized (right).
    Left: single central learner with full joint state.
    Centre-left: CTDE --- centralized critic, decentralized actors.
    Centre-right: networked agents --- local gradients + neighborhood consensus.
    Right: independent Q-learning --- no communication during training.
    Arrows indicate the direction of increasing deployability vs.\ policy quality.
  }}
  \caption{Spectrum of training paradigms in cooperative \gls{MARL}.}
  \label{fig:decentralized-spectrum}
\end{figure}
